(1) 平面波的平均总电磁能量密度为 \(u=I/c\)，动量密度为 \(\boldsymbol g=\boldsymbol S/c^2=(I/c^2)\hat z\)。在真空中能量以 c 传播，故单位面积每秒携带的能量是 I、动量是 I/c。

(2) 全吸收板每单位时间吸收动量大小 \(IA/c\)，板受到的平均法向力 \(F_{\rm abs}=IA/c\)，辐射压
\[
p_{\rm abs}=F_{\rm abs}/A=I/c.
\]
能量沉积为板的热能或其他内部能，这不影响入射动量吸收率。

(3) 理想反射镜把每份入射光子沿 z 的动量从 \(+p\) 改为 \(-p\)，光的动量改变为 \(-2p\)，镜获得 \(+2p\)。故
\[
F_{\rm ref}=2IA/c,\qquad p_{\rm ref}=2I/c.
\]
系数2并不要求光在反射时损失额外能量：静止理想镜使光子能量保持不变，但方向反转。

(4) 取每个光子能量 \(\hbar\omega\)，动量 \(\hbar\omega/c\)，光子流率 \(\dot N=IA/(\hbar\omega)\)。全吸收传递动量 \(\dot N\hbar\omega/c=IA/c\)；全反射传递两倍。与经典 Maxwell 应力动量流结果一致。
